
\documentclass[11pt,a4paper]{article}
\usepackage[utf8]{inputenc}
\usepackage[T1]{fontenc}
\usepackage{amsmath}
\usepackage{geometry}
\usepackage{xcolor}
\usepackage{listings}

\geometry{margin=1in}

\title{\textbf{UAV-Guardian: Edge-AI Anomaly Detection via Deep Autoencoders}}
\author{AEC Final Year Project}
\date{\today}

\begin{document}
\maketitle

\section{Introduction}
UAV-Guardian is a real-time predictive maintenance system designed for Unmanned Aerial Vehicles (UAVs). It utilizes an \textbf{Edge-deployed Deep Autoencoder} on an ESP32 to detect mechanical anomalies by analyzing high-frequency vibrational data from an MPU6050 IMU.

\section{The AI Engine: Unsupervised Anomaly Detection}
This project employs an \textbf{Unsupervised Autoencoder}. This approach is critical for drone safety as failure modes are diverse and often unpredictable.

\subsection{Architecture}
\begin{itemize}
    \item \textbf{Input Layer:} 300 features (1s of X, Y, Z data at 100Hz).
    \item \textbf{Encoder:} 300 $\rightarrow$ 32 $\rightarrow$ 8 (Latent Space).
    \item \textbf{Decoder:} 8 $\rightarrow$ 32 $\rightarrow$ 300 (Reconstruction).
\end{itemize}

\subsection{Mathematical Framework}
Anomaly detection is measured via Mean Squared Error (MSE):
\begin{equation}
MSE = \frac{1}{n} \sum_{i=1}^{n} (x_i - \hat{x}_i)^2
\end{equation}

\section{Edge Deployment}
The model was converted to \texttt{.tflite} and embedded into \texttt{uav_model.h}. A manual input backup buffer was used to resolve the TFLite memory address conflict.

\end{document}
